% Unit 051 terjemahan Bahasa Indonesia dari chapter4.tex baris 260--560.
% Sumber: Methods of Algebra, Volume 2, commit
% 9a5803ff77dd3257484cb177f851a73770a59dd3 (CC BY 4.0).
% stable-unit-id: o014.aljabr2.chapter4.triangulated-basics
% Cakupan: Bagian 4.2 mengenai sifat dasar kategori bertriangulasi; berhenti
% sebelum Bagian 4.3 pada baris 561.

% segment-id: o014.li2.chapter4.triangulated-basics.h001
\section{Sifat Dasar}\label{sec:triangular-basic}
\hypertarget{o014.aljabr2.chapter4.triangulated-basics}{ }

% segment-id: o014.li2.chapter4.triangulated-basics.p001
Bagian sebelumnya hanya menyajikan definisi dan sifat awal kategori
bertriangulasi dan funktor kohomologis. Bagian ini menguraikan struktur lebih
lanjut.

% segment-id: o014.li2.chapter4.triangulated-basics.pr001
\begin{proposisi}\label{prop:dt-3-2}
	Tinjau diagram komutatif berikut dalam suatu kategori prabertriangulasi
	$\mathcal{D}$:
% segment-id: o014.li2.chapter4.triangulated-basics.g001
	\[\begin{tikzcd}
		X \arrow[r] \arrow[d, "\alpha"'] & Y \arrow[d, "\beta"] \arrow[r] & Z \arrow[d, "\gamma"] \arrow[r] & TX \arrow[d, "T\alpha"] \\
		X' \arrow[r] & Y' \arrow[r] & Z' \arrow[r] & TX'
	\end{tikzcd}\]
	Jika setiap baris merupakan segitiga terbedakan dan dua mana pun di antara
	$\alpha$, $\beta$, dan $\gamma$ merupakan isomorfisme, maka morfisme yang
	tersisa juga merupakan isomorfisme.
\end{proposisi}

% segment-id: o014.li2.chapter4.triangulated-basics.pf001
\begin{bukti}
	Setelah menerapkan rotasi (TR3) yang sesuai, tanpa mengurangi keumuman kita
	dapat mengandaikan bahwa $\alpha$ dan $\gamma$ merupakan isomorfisme. Untuk
	membuktikan bahwa $\beta$ merupakan isomorfisme, cukup dibuktikan bagi setiap
	objek $S$ bahwa
	$\Hom(S, Y) \xrightarrow{\beta_*} \Hom(S, Y')$ merupakan isomorfisme.
	Proposisi~\ref{prop:cohomological-Hom-functor} memberikan diagram komutatif
	dengan baris-baris eksak berikut:
% segment-id: o014.li2.chapter4.triangulated-basics.g002
	\[\begin{tikzcd}[column sep=small]
		\Hom(S, T^{-1} Z) \arrow[d, "\sim" sloped] \arrow[r] & \Hom(S, X) \arrow[r] \arrow[d, "\sim" sloped] & \Hom(S,Y) \arrow[d, "{\beta_*}"] \arrow[r] & \Hom(S,Z) \arrow[d, "\sim" sloped] \arrow[r] & \Hom(S, TX) \arrow[d, "\sim" sloped] \\
		\Hom(S, T^{-1} Z') \arrow[r] & \Hom(S, X') \arrow[r] & \Hom(S, Y') \arrow[r] & \Hom(S, Z') \arrow[r] & \Hom(S, TX') .
	\end{tikzcd}\]
	Sekarang terapkan Proposisi~\ref{prop:5-lemma} pada diagram tersebut.
\end{bukti}

% segment-id: o014.li2.chapter4.triangulated-basics.p002
Dalam argumen di atas, $\Hom(S, \cdot)$ juga dapat diganti dengan
$\Hom(\cdot, S)$.

% segment-id: o014.li2.chapter4.triangulated-basics.r001
\begin{catatan}\label{rem:special-triangles}
	Untuk sementara, sebut segitiga $X \to Y \to Z \xrightarrow{+1}$ sebagai
	\emph{segitiga khusus} (atau \emph{segitiga kokhusus}) apabila, untuk setiap
	$S \in \Obj(\mathcal{D})$, barisan
% segment-id: o014.li2.chapter4.triangulated-basics.q001
	\[ \cdots \to \Hom(S, X) \to \Hom(S, Y) \to \Hom(S, Z) \to \Hom(S, TX) \to \cdots \]
	(atau versi yang bersesuaian dengan $\Hom(\cdot, S)$) eksak. Jika
	$X_i \to Y_i \to Z_i \xrightarrow{+1}$ merupakan suatu keluarga segitiga
	terbedakan ($i \in I$), maka
	$\prod_i X_i \to \prod_i Y_i \to \prod_i Z_i \xrightarrow{+1}$ (atau
	$\coprod_i X_i \to \coprod_i Y_i \to \coprod_i Z_i \xrightarrow{+1}$)
	merupakan segitiga khusus (atau segitiga kokhusus), asalkan produk (atau
	koproduk) yang dibicarakan ada. Hal ini semata-mata mengikuti sifat universal
	produk dan koproduk bersama hasil \cite[Lema~6.8.3]{Li1}. Perhatikan bahwa,
	karena $T: \mathcal{D} \to \mathcal{D}$ merupakan suatu ekuivalensi, $T$
	mesti mempertahankan produk dan koproduk. Proposisi~\ref{prop:dt-3-2} dapat
	diperluas ke kasus ketika kedua baris merupakan segitiga khusus (atau
	segitiga kokhusus).
\end{catatan}

% segment-id: o014.li2.chapter4.triangulated-basics.co001
\begin{korolari}\label{prop:Cone-uniqueness}
	Untuk setiap morfisme $f: X \to Y$ dalam suatu kategori prabertriangulasi
	$\mathcal{D}$, segitiga terbedakan
	$X \xrightarrow{f} Y \to Z \xrightarrow{+1}$ dalam aksioma (TR2) unik
	hingga isomorfisme.
\end{korolari}

% segment-id: o014.li2.chapter4.triangulated-basics.pf002
\begin{bukti}
	Misalkan diberikan segitiga-segitiga terbedakan
	$X \xrightarrow{f} Y \to Z \xrightarrow{+1}$ dan
	$X \xrightarrow{f} Y \to Z' \xrightarrow{+1}$, aksioma (TR4) memberikan
	morfisme $\gamma$ yang membuat diagram
% segment-id: o014.li2.chapter4.triangulated-basics.g003
	\[\begin{tikzcd}
		X \arrow[r, "f"] \arrow[d, "\identity_X"'] & Y \arrow[r] \arrow[d, "\identity_Y"] & Z \arrow[r] \arrow[d, "\gamma"] & TX \arrow[d, "\identity_{TX}"] \\
		X \arrow[r, "f"'] & Y \arrow[r] & Z' \arrow[r] & T X
	\end{tikzcd}\]
	komutatif. Proposisi~\ref{prop:dt-3-2} kemudian menyiratkan bahwa $\gamma$
	merupakan isomorfisme.
\end{bukti}

% segment-id: o014.li2.chapter4.triangulated-basics.p003
Perlu ditekankan bahwa isomorfisme dalam
Korolari~\ref{prop:Cone-uniqueness} \emph{tidak kanonik}.

% segment-id: o014.li2.chapter4.triangulated-basics.le001
\begin{lema}\label{prop:dt-sum}
	Misalkan $\mathcal{D}$ kategori prabertriangulasi dan
	$X_i \xrightarrow{f_i} Y_i \xrightarrow{g_i} Z_i \xrightarrow{h_i} TX_i$
	suatu keluarga segitiga terbedakan. Andaikan $\prod_i X_i$,
	$\prod_i Y_i$, dan $\prod_i Z_i$ semuanya ada. Maka
% segment-id: o014.li2.chapter4.triangulated-basics.q002
	\[ \prod_{i \in I} X_i \xrightarrow{\prod_i f_i } \prod_{i \in I} Y_i \xrightarrow{\prod_i g_i} \prod_{i \in I} Z_i \xrightarrow{\prod_i h_i} \underbracket{\prod_{i \in I} TX_i}_{\simeq T(\prod_{i \in I} X_i)} \]
	juga merupakan segitiga terbedakan. Pernyataan yang bersesuaian tetap
	berlaku jika $\prod_{i \in I}$ diganti dengan $\coprod_{i \in I}$.
\end{lema}

% segment-id: o014.li2.chapter4.triangulated-basics.pf003
\begin{bukti}
	Berdasarkan dualitas, kita hanya menangani kasus $\prod_{i \in I}$. Mulai
	sekarang, kita selalu mengidentifikasi $\prod_i TX_i$ dengan
	$T(\prod_i X_i)$. Dengan (TR2), pilih segitiga terbedakan
	$\prod_i X_i \xrightarrow{\prod_i f_i} \prod_i Y_i \to Q
	\xrightarrow{h} \prod_i TX_i$. Untuk setiap $i \in I$, terapkan (TR4)
	untuk memperoleh diagram komutatif
% segment-id: o014.li2.chapter4.triangulated-basics.g004
	\[\begin{tikzcd}
		\prod_{j \in I} X_j \arrow[r, "{\prod_j f_j}"] \arrow[d] & \prod_j Y_j \arrow[r] \arrow[d] & Q \arrow[r, "h"] \arrow[d, "{\exists \gamma_i}"] & \prod_{j \in I} TX_j \arrow[d] \\
		X_i \arrow[r, "{f_i}"'] & Y_i \arrow[r, "{g_i}"'] & Z_i \arrow[r, "{h_i}"'] & TX_i
	\end{tikzcd}\]
	Semua morfisme vertikal selain $\gamma_i$ merupakan morfisme proyeksi
	kanonik. Sifat universal produk kemudian memberikan diagram komutatif
% segment-id: o014.li2.chapter4.triangulated-basics.g005
	\[\begin{tikzcd}
		\prod_{i \in I} X_i \arrow[r, "{\prod_i f_i}"] \arrow[d, "\identity"'] & \prod_{i \in I} Y_i \arrow[r] \arrow[d, "\identity"] & Q \arrow[r, "h"] \arrow[d, "{\prod_i \gamma_i}"] & \prod_{i \in I} TX_i \arrow[d, "\identity"] \\
		\prod_{i \in I} X_i \arrow[r, "{\prod_i f_i}"'] & \prod_{i \in I} Y_i \arrow[r, "{\prod_i g_i}"'] & \prod_{i \in I} Z_i \arrow[r, "{\prod_i h_i}"'] & \prod_{i \in I} TX_i .
	\end{tikzcd}\]
	Baris kedua merupakan segitiga khusus dalam pengertian
	Catatan~\ref{rem:special-triangles}. Perluasan
	Proposisi~\ref{prop:dt-3-2} kemudian menyiratkan bahwa
	$\prod_{i \in I} \gamma_i$ merupakan isomorfisme.
\end{bukti}

% segment-id: o014.li2.chapter4.triangulated-basics.co002
\begin{korolari}\label{prop:dt-sum-0}
	Dalam setiap kategori prabertriangulasi, segitiga berbentuk
	$X \xrightarrow{\iota_1} X \oplus Y \xrightarrow{p_2} Y \xrightarrow{0} TX$
	selalu merupakan segitiga terbedakan; di sini $\iota_i$ dan $p_j$
	didefinisikan seperti dalam \S\ref{sec:Ker-Coker}.
\end{korolari}

% segment-id: o014.li2.chapter4.triangulated-basics.pf004
\begin{bukti}
	Terapkan Lema~\ref{prop:dt-sum} pada segitiga-segitiga terbedakan
	$X \xrightarrow{\identity_X} X \to 0 \to TX$ dan
	$0 \to Y \xrightarrow{\identity_Y} Y \to 0$.
\end{bukti}

% segment-id: o014.li2.chapter4.triangulated-basics.p004
Sebaliknya, jika salah satu morfisme dalam suatu segitiga terbedakan adalah
$0$, segitiga itu berasal dari suatu jumlah langsung. Fakta ini dapat
diverifikasi dengan merotasinya ke bentuk berikut.

% segment-id: o014.li2.chapter4.triangulated-basics.pr002
\begin{proposisi}\label{prop:dt-sum-1}
	Dalam setiap kategori prabertriangulasi, suatu segitiga terbedakan berbentuk
	$X \to M \to Y \xrightarrow{0} TX$ mesti isomorfik dengan segitiga
	$X \to X \oplus Y \to Y \xrightarrow{0} TX$ dalam
	Korolari~\ref{prop:dt-sum-0}.
\end{proposisi}

% segment-id: o014.li2.chapter4.triangulated-basics.pf005
\begin{bukti}
	Versi (TR4) yang telah dirotasi secara tepat menunjukkan bahwa terdapat
	morfisme $\alpha: X \oplus Y \to M$ sedemikian sehingga diagram berikut
	merupakan morfisme di antara segitiga-segitiga terbedakan:
% segment-id: o014.li2.chapter4.triangulated-basics.g006
	\[\begin{tikzcd}
		X \arrow[r, "\iota_1"] \arrow[d, "\identity_X"'] & X \oplus Y \arrow[r, "p_2"] \arrow[d, "\alpha"] & Y \arrow[r, "0"] \arrow[d, "\identity_Y"] & TX \arrow[d, "\identity_{TX}"] \\
		X \arrow[r] & M  \arrow[r] & Y \arrow[r, "0"']  & TX \\
	\end{tikzcd}\]
	Proposisi~\ref{prop:dt-3-2} menyiratkan bahwa $\alpha$ merupakan
	isomorfisme.
\end{bukti}

% segment-id: o014.li2.chapter4.triangulated-basics.p005
Hasil berikut menunjukkan bahwa aditivitas dalam definisi funktor
bertriangulasi sebenarnya berlebihan.

% segment-id: o014.li2.chapter4.triangulated-basics.co003
\begin{korolari}\label{prop:triangulated-automatic-additivity}
	Misalkan $\mathcal{D}$ dan $\mathcal{D}'$ kategori-kategori
	prabertriangulasi yang funktor pergeserannya sama-sama dinotasikan dengan
	$T$. Misalkan $F: (\mathcal{D}, T) \to (\mathcal{D}', T)$ suatu funktor di
	antara kategori-kategori dengan pergeseran. Jika $F$ memetakan segitiga
	terbedakan ke segitiga terbedakan, maka $F$ aditif dan oleh karena itu secara
	otomatis merupakan funktor bertriangulasi.
\end{korolari}

% segment-id: o014.li2.chapter4.triangulated-basics.pf006
\begin{bukti}
	Segitiga terbedakan
	$0 \xrightarrow{\identity} 0 \xrightarrow{\identity} 0
	\xrightarrow{\identity} 0$ di $\mathcal{D}$ memberikan segitiga
	terbedakan
	$F(0) \xrightarrow{\identity} F(0) \xrightarrow{\identity} F(0)
	\xrightarrow{\identity} F(0)$ di $\mathcal{D}'$.
	Lema~\ref{prop:dt-composite-null} kemudian memberikan
	$\identity_{F(0)} = \identity_{F(0)} \circ \identity_{F(0)} = 0$, sehingga
	$F(0) \simeq 0$. Hal ini juga menyiratkan bahwa $F$ memetakan morfisme nol
	ke morfisme nol.

	Selanjutnya, tinjau segitiga terbedakan di $\mathcal{D}$ yang diberikan
	oleh Korolari~\ref{prop:dt-sum-0}, yakni
	$X \xrightarrow{\iota_1} X \oplus Y \xrightarrow{p_2} Y
	\xrightarrow{0} TX$. Segitiga yang bersesuaian
% segment-id: o014.li2.chapter4.triangulated-basics.q003
	\[ FX \xrightarrow{F\iota_1} F(X \oplus Y) \xrightarrow{Fp_2} FY \xrightarrow{0} TFX \]
	merupakan segitiga terbedakan di $\mathcal{D}'$. Ambil
	$M := F(X \oplus Y)$ dalam Proposisi~\ref{prop:dt-sum-1}. Morfisme $\alpha$
	dalam bukti proposisi tersebut jelas dapat diambil sebagai morfisme kanonik
	$FX \oplus FY \to F(X \oplus Y)$ yang ditentukan oleh sifat universal
	jumlah langsung. Jadi, $\alpha$ merupakan isomorfisme. Berdasarkan
	Korolari~\ref{prop:automatic-additivity} (iii), $F$ aditif.
\end{bukti}

% segment-id: o014.li2.chapter4.triangulated-basics.p006
Dua hasil berikut menunjukkan bahwa, dalam suatu ekuivalensi kategori
prabertriangulasi (Definisi~\ref{def:triangulated-cat}), cukup diasumsikan bahwa
salah satu arahnya merupakan funktor bertriangulasi. Ini adalah penerapan lain
Proposisi~\ref{prop:cohomological-Hom-functor}.

% segment-id: o014.li2.chapter4.triangulated-basics.pr003
\begin{proposisi}\label{prop:adjoint-triangulated}
	Misalkan $\mathcal{D}$ dan $\mathcal{D}'$ kategori-kategori
	prabertriangulasi, dan tinjau funktor-funktor
% segment-id: o014.li2.chapter4.triangulated-basics.g007
	$\begin{tikzcd}
		\mathcal{D} \arrow[shift left, r, "F"] & \mathcal{D}' \arrow[shift left, l, "G"]
	\end{tikzcd}$
	dengan $(F, G)$ suatu adjungsi. Maka $F$ merupakan funktor
	bertriangulasi jika dan hanya jika $G$ merupakan funktor bertriangulasi.
	Dalam keadaan ini, unit $\eta: \identity_{\mathcal{D}} \to GF$ dan kounit
	$\varepsilon: FG \to \identity_{\mathcal{D}'}$ dari adjungsi tersebut
	kompatibel dengan funktor pergeseran; lihat
	Definisi~\ref{def:category-with-translation}.
\end{proposisi}

% segment-id: o014.li2.chapter4.triangulated-basics.pf007
\begin{bukti}
	Berdasarkan dualitas, cukup ditinjau kasus ketika $F$ merupakan funktor
	bertriangulasi. Korolari~\ref{prop:automatic-additivity} menunjukkan bahwa
	$G$ aditif. Tuliskan kedua funktor pergeseran pada $\mathcal{D}$ dan
	$\mathcal{D}'$ dengan $T$. Dengan memindahkan struktur, kita memperoleh
	adjungsi baru
% segment-id: o014.li2.chapter4.triangulated-basics.q004
	\[ \hat{F} := T F T^{-1}, \quad \hat{G} := T G T^{-1}, \quad \hat{\eta} := T\eta T^{-1}, \quad \hat{\varepsilon} := T \varepsilon T^{-1}. \]
	Karena $\hat{F} \simeq F$, keunikan adjoin kanan
	\cite[Proposisi~2.6.10]{Li1} memberikan isomorfisme yang bersesuaian
	$\hat{G} \simeq G$ beserta diagram-diagram komutatif
% segment-id: o014.li2.chapter4.triangulated-basics.g008
	\[\begin{tikzcd}
		\hat{F} \hat{G} \arrow[r, "\hat{\varepsilon}"] \arrow[d, "\sim" sloped] & \identity_{\mathcal{D}'} \\
		FG \arrow[ru, "{\varepsilon}"'] &
	\end{tikzcd} \quad \begin{tikzcd}
		\hat{G} \hat{F} \arrow[d, "\sim" sloped] & \identity_{\mathcal{D}} \arrow[l, "\hat{\eta}"'] \arrow[ld, "\eta"] \\
		GF &
	\end{tikzcd}\]
	Secara khusus, $GT \simeq TG$ dan $FGT \simeq FTG \simeq TFG$.
	Setelah diagram komutatif untuk $\hat{\varepsilon}$ dan $\varepsilon$
	dikomposisikan di kanan dengan $T$, diagram tersebut dapat ditulis ulang
	sebagai
% segment-id: o014.li2.chapter4.triangulated-basics.g009
	\begin{equation}\label{eqn:adjoint-triangulated-aux0}\begin{tikzcd}
		TFG \arrow[r, "{T\varepsilon}"] \arrow[d, "\sim" sloped] & T \\
		FGT \arrow[ru, "\varepsilon T"'] &
	\end{tikzcd}\end{equation}
	Persamaan ini dan versi dualnya untuk $\eta$ menunjukkan bahwa $\eta$ dan
	$\varepsilon$ kompatibel dengan funktor pergeseran.

	Sekarang kita tunjukkan bahwa $G$ mempertahankan segitiga terbedakan. Tinjau
	segitiga terbedakan $X \xrightarrow{f} Y \xrightarrow{g} Z
	\xrightarrow{+1}$ di $\mathcal{D}'$. Pilih segitiga terbedakan
	$GX \xrightarrow{Gf} GY \xrightarrow{h} A \xrightarrow{+1}$ di
	$\mathcal{D}$, terapkan $F$ kepadanya, lalu gunakan (TR3) untuk
	memperoleh diagram komutatif
% segment-id: o014.li2.chapter4.triangulated-basics.g010
	\begin{equation*}\begin{tikzcd}
		FG X \arrow[r, "FGf"] \arrow[d, "\varepsilon_X"'] & FG Y \arrow[d, "\varepsilon_Y"] \arrow[r, "Fh"] & FA \arrow[r] \arrow[d] & TFGX \arrow[d, "T\varepsilon_X"] \\
		X \arrow[r, "f"'] & Y \arrow[r, "g"'] & Z \arrow[r] & TX .
	\end{tikzcd}\end{equation*}

	Terapkan adjungsi pada diagram tersebut, lalu sesuaikan persegi paling kanan
	dengan \eqref{eqn:adjoint-triangulated-aux0} dan $GT \simeq TG$. Dengan
	mudah diperoleh diagram komutatif
% segment-id: o014.li2.chapter4.triangulated-basics.g011
	\[\begin{tikzcd}
		GX \arrow[r, "Gf"] \arrow[d, "\identity_X"'] & GY \arrow[d, "\identity_Y"] \arrow[r, "h"] & A \arrow[r] \arrow[d] & TGX \arrow[d, "\identity_{TGX}"] \\
		GX \arrow[r, "Gf"'] & GY \arrow[r, "Gg"'] & GZ \arrow[r] & TGX
	\end{tikzcd}\]
	Baris pertama merupakan segitiga terbedakan. Untuk setiap
	$B \in \Obj(\mathcal{D})$, terapkan $\Hom_{\mathcal{D}}(B, \cdot)$ untuk
	memperoleh diagram komutatif
% segment-id: o014.li2.chapter4.triangulated-basics.g012
	\[\begin{tikzcd}[column sep=small, every cell/.append style={font=\small}]
		\Hom(B, GX) \arrow[r, "(Gf)_*" inner sep=0.8em] \arrow[d, "\simeq"'] & \Hom(B, GY) \arrow[r, "h_*" inner sep=0.8em] \arrow[d, "\simeq"] & \Hom(B, A) \arrow[d] \arrow[r] & \Hom(B, TGX) \arrow[r] \arrow[d, "\simeq"] & \Hom(B, TGY) \arrow[d, "\simeq"] \\
		\Hom(B, GX) \arrow[r, "(Gf)_*"' inner sep=0.8em] & \Hom(B, GY) \arrow[r, "(Gg)_*"' inner sep=0.8em] & \Hom(B, GZ) \arrow[r] & \Hom(B, TGX) \arrow[r] & \Hom(B, TGY)
	\end{tikzcd}\]
	Proposisi~\ref{prop:cohomological-Hom-functor} menunjukkan bahwa baris
	pertama eksak. Di sisi lain, berdasarkan $TG \simeq GT$ dan adjungsi, baris
	kedua isomorfik dengan barisan yang diinduksi oleh segitiga terbedakan
	$X \xrightarrow{f} Y \xrightarrow{g} Z \xrightarrow{+1}$, yakni
% segment-id: o014.li2.chapter4.triangulated-basics.g013
	\[\begin{tikzcd}[column sep=small, every cell/.append style = {font = \small}]
		\Hom(FB, X) \arrow[r, "f_*"] & \Hom(FB, Y) \arrow[r, "g_*"] & \Hom(FB, Z) \arrow[r] & \Hom(FB, TX) \arrow[r] & \Hom(FB, TY) ,
	\end{tikzcd}\]
	dan karena itu juga eksak. Karena $B$ sembarang,
	Proposisi~\ref{prop:5-lemma} menyiratkan bahwa $A \to GZ$ merupakan
	isomorfisme. Jadi,
	$GX \xrightarrow{Gf} GY \xrightarrow{Gg} GZ \to TGX$ memang merupakan
	segitiga terbedakan di $\mathcal{D}$.
\end{bukti}

% segment-id: o014.li2.chapter4.triangulated-basics.co004
\begin{korolari}\label{prop:triangulated-equivalence}
	Misalkan $\mathcal{D}$ dan $\mathcal{D}'$ kategori-kategori
	prabertriangulasi dan funktor-funktor
% segment-id: o014.li2.chapter4.triangulated-basics.g014
	$\begin{tikzcd}
		\mathcal{D} \arrow[shift left, r, "F"] & \mathcal{D}' \arrow[shift left, l, "G"]
	\end{tikzcd}$
	saling kuasi-invers. Maka:
% segment-id: o014.li2.chapter4.triangulated-basics.l001
	\begin{compactenum}[(i)]
% segment-id: o014.li2.chapter4.triangulated-basics.i001
		\item $F$ merupakan funktor bertriangulasi jika dan hanya jika $G$ juga
		merupakan funktor bertriangulasi;
% segment-id: o014.li2.chapter4.triangulated-basics.i002
		\item jika $F$ merupakan funktor bertriangulasi, maka $F$ dan $G$
		memberikan ekuivalensi antara kategori-kategori prabertriangulasi.
	\end{compactenum}
\end{korolari}

% segment-id: o014.li2.chapter4.triangulated-basics.pf008
\begin{bukti}
	Gunakan \cite[Teorema~2.6.12]{Li1} untuk merealisasikan $(F, G)$ dan
	$(G, F)$ masing-masing sebagai adjungsi, lalu terapkan
	Proposisi~\ref{prop:adjoint-triangulated}.
\end{bukti}

% segment-id: o014.li2.chapter4.triangulated-basics.dp001
\begin{definisiproposisi}[Subkategori Prabertriangulasi]\label{def:triangulated-subcat}
	\index{kategori prabertriangulasi!subkategori prabertriangulasi}
	\index{kategori bertriangulasi!subkategori bertriangulasi}
	Misalkan $\mathcal{D}'$ subkategori penuh aditif dari suatu kategori
	prabertriangulasi $\mathcal{D}$ yang tertutup terhadap funktor pergeseran
	$T$. Nyatakan funktor inklusi dengan
	$\iota: \mathcal{D}' \to \mathcal{D}$.
% segment-id: o014.li2.chapter4.triangulated-basics.l002
	\begin{enumerate}[(i)]
% segment-id: o014.li2.chapter4.triangulated-basics.i003
		\item Jika ada, struktur prabertriangulasi pada $\mathcal{D}'$ yang
		membuat $\iota$ menjadi funktor bertriangulasi bersifat unik. Dalam
		keadaan ini, suatu segitiga $X \to Y \to Z \xrightarrow{+1}$ di
		$\mathcal{D}'$ merupakan segitiga terbedakan jika dan hanya jika ia
		merupakan segitiga terbedakan di $\mathcal{D}$.
% segment-id: o014.li2.chapter4.triangulated-basics.i004
		\item Struktur prabertriangulasi tersebut ada jika dan hanya jika syarat
		berikut terpenuhi: apabila $X \to Y \to Z \xrightarrow{+1}$ merupakan
		segitiga terbedakan di $\mathcal{D}$ dan sembarang dua sukunya termasuk
		dalam $\Obj(\mathcal{D}')$, maka suku yang lain isomorfik dengan suatu
		objek di $\mathcal{D}'$.
% segment-id: o014.li2.chapter4.triangulated-basics.i005
		\item Jika lebih lanjut $\mathcal{D}$ diasumsikan sebagai kategori
		bertriangulasi, maka $\mathcal{D}'$ juga merupakan kategori
		bertriangulasi.
	\end{enumerate}

	Kategori prabertriangulasi (atau kategori bertriangulasi) $\mathcal{D}'$
	semacam ini disebut \emph{subkategori prabertriangulasi} (atau
	\emph{subkategori bertriangulasi}) dari $\mathcal{D}$.
\end{definisiproposisi}

% segment-id: o014.li2.chapter4.triangulated-basics.pf009
\begin{bukti}
	Pertama, kita tangani (i). Andaikan $\iota$ merupakan funktor
	bertriangulasi. Tentu saja, setiap segitiga terbedakan di $\mathcal{D}'$
	merupakan segitiga terbedakan di $\mathcal{D}$. Sebaliknya, andaikan
	segitiga terbedakan $X \xrightarrow{f} Y \to Z \xrightarrow{+1}$ di
	$\mathcal{D}$ memenuhi $X, Y, Z \in \Obj(\mathcal{D}')$. Pilih sembarang
	segitiga terbedakan $X \xrightarrow{f} Y \to Z' \xrightarrow{+1}$ di
	$\mathcal{D}'$. Dengan menerapkan Korolari~\ref{prop:Cone-uniqueness} di
	$\mathcal{D}$, kedua segitiga tersebut isomorfik. Karena
	$\mathcal{D}'$ merupakan subkategori penuh, (TR0) menyiratkan bahwa
	$X \xrightarrow{f} Y \to Z \xrightarrow{+1}$ juga merupakan segitiga
	terbedakan di $\mathcal{D}'$. Inilah tepatnya struktur prabertriangulasi yang
	dinyatakan dalam (i).

	Jika struktur prabertriangulasi tersebut ada, maka
	Korolari~\ref{prop:Cone-uniqueness} yang telah dibuktikan menunjukkan bahwa
	syarat dalam (ii) perlu, sebab setelah rotasi kita selalu dapat mengandaikan
	$X, Y \in \Obj(\mathcal{D}')$. Sebaliknya, andaikan syarat (ii) terpenuhi,
	dan definisikan segitiga terbedakan di $\mathcal{D}'$ dengan cara yang
	dinyatakan dalam (i). Aksioma (TR0), (TR1), (TR3), dan (TR4) untuk
	$\mathcal{D}'$ diwarisi langsung dari $\mathcal{D}$. Untuk (TR2), bagi suatu
	morfisme $f: X \to Y$ di $\mathcal{D}'$, pilih segitiga terbedakan
	$X \xrightarrow{f} Y \to Z \xrightarrow{+1}$ di $\mathcal{D}$. Syarat
	tersebut menyatakan bahwa $Z$ isomorfik dengan suatu objek di
	$\mathcal{D}'$; jadi, kita boleh mengandaikan bahwa
	$Z \in \Obj(\mathcal{D}')$, dan dengan demikian memperoleh segitiga
	terbedakan di $\mathcal{D}'$. Terbukti bahwa $\mathcal{D}'$ menjadi kategori
	prabertriangulasi. Ini membuktikan kecukupan.

	Untuk (iii), intinya ialah memverifikasi (TR5) bagi $\mathcal{D}'$.
	Sesungguhnya, segitiga $Z' \to Y' \to X' \xrightarrow{+1}$ yang diperoleh
	dengan menerapkan (TR5) di $\mathcal{D}$ juga secara otomatis merupakan
	segitiga terbedakan di $\mathcal{D}'$ menurut (i).
\end{bukti}

% segment-id: o014.li2.chapter4.triangulated-basics.cv001
\begin{konvensi}\label{con:saturated-intersection}
	\index{subkategori!jenuh (replete/saturated subcategory)}
	Tetapkan suatu kategori $\mathcal{C}$. Suatu subkategori $\mathcal{C}'$
	disebut \emph{subkategori jenuh} apabila tertutup terhadap isomorfisme,
	yakni
% segment-id: o014.li2.chapter4.triangulated-basics.q005
	\[ \forall X \in \Obj(\mathcal{C}'), \; \forall Y \in \Obj(\mathcal{C}), \quad X \simeq Y \implies Y \in \Obj(\mathcal{C}') . \]
	Untuk setiap keluarga subkategori penuh $(\mathcal{C}_i)_{i \in I}$, dengan
	$I$ tidak kosong, irisannya $\bigcap_{i \in I} \mathcal{C}_i$ menurut
	definisi merupakan subkategori penuh dari $\mathcal{C}$ yang memenuhi
	$\Obj\left( \bigcap_{i \in I} \mathcal{C}_i \right)
	= \bigcap_{i \in I} \Obj(\mathcal{C}_i)$.
\end{konvensi}

% segment-id: o014.li2.chapter4.triangulated-basics.pr004
\begin{proposisi}\label{prop:triangulated-subcat-intersection}
	Misalkan $\mathcal{D}$ kategori prabertriangulasi (atau kategori
	bertriangulasi), dan $(\mathcal{D}_i)_{i \in I}$ suatu keluarga subkategori
	prabertriangulasi jenuh (atau subkategori bertriangulasi jenuh), dengan
	$I \neq \emptyset$. Maka $\bigcap_{i \in I} \mathcal{D}_i$ juga merupakan
	subkategori prabertriangulasi jenuh (atau subkategori bertriangulasi jenuh).
\end{proposisi}

% segment-id: o014.li2.chapter4.triangulated-basics.pf010
\begin{bukti}
	Irisan subkategori-subkategori jenuh jelas mempertahankan syarat dalam
	Definisi--Proposisi~\ref{def:triangulated-subcat} (ii).
\end{bukti}

% segment-id: o014.li2.chapter4.triangulated-basics.p007
Suatu teknik yang umum ialah memotong subkategori prabertriangulasi jenuh
(atau subkategori bertriangulasi jenuh) dengan menggunakan funktor
kohomologis.

% segment-id: o014.li2.chapter4.triangulated-basics.pr005
\begin{proposisi}\label{prop:cut-subcat-by-cohomologies}
	Misalkan $\mathcal{D}$ kategori prabertriangulasi (atau kategori
	bertriangulasi), $H: \mathcal{D} \to \mathcal{A}$ suatu funktor
	kohomologis, dan $\mathcal{T}$ suatu subkategori Serre lemah dari
	$\mathcal{A}$ (Definisi~\ref{def:Serre-subcat}). Definisikan subkategori
	penuh $\mathcal{D}_{H, \mathcal{T}}$ dari $\mathcal{D}$ dengan
% segment-id: o014.li2.chapter4.triangulated-basics.q006
	\[ X \in \Obj(\mathcal{D}_{H, \mathcal{T}}) \iff \forall n \in \Z, \; H(T^n X) \in \Obj(\mathcal{T}). \]
	Maka $\mathcal{D}_{H, \mathcal{T}}$ merupakan subkategori
	prabertriangulasi jenuh (atau subkategori bertriangulasi jenuh) dari
	$\mathcal{D}$.
\end{proposisi}

% segment-id: o014.li2.chapter4.triangulated-basics.pf011
\begin{bukti}
	Cukup memverifikasi syarat dalam
	Definisi--Proposisi~\ref{def:triangulated-subcat} (ii). Ingat bahwa suatu
	subkategori Serre lemah bersifat jenuh. Karena itu, sifat jenuh dan
	$T$-invariansi dari $\mathcal{D}_{H, \mathcal{T}}$ sudah jelas. Sekarang
	tinjau segitiga terbedakan $X \to Y \to Z \xrightarrow{+1}$ di
	$\mathcal{D}$ yang dua sukunya berada dalam
	$\mathcal{D}_{H, \mathcal{T}}$. Setelah merotasi, tanpa mengurangi keumuman
	kita dapat mengandaikan bahwa
	$X, Z \in \Obj(\mathcal{D}_{H, \mathcal{T}})$. Untuk setiap $n \in \Z$,
	terdapat barisan eksak
% segment-id: o014.li2.chapter4.triangulated-basics.q007
	\[ H(T^{n-1} Z) \to H(T^n X) \to H(T^n Y) \to H(T^n Z) \to H(T^{n+1} X). \]
	Dengan menerapkan definisi subkategori Serre lemah, segera diperoleh
	$H(T^n Y) \in \Obj(\mathcal{T})$.
\end{bukti}

% segment-id: o014.li2.chapter4.triangulated-basics.ex001
\begin{contoh}\label{eg:subcat-N-H}
	\index[sym1]{NH@$\mathcal{N}_H$}
	Misalkan $\mathcal{D}$ kategori prabertriangulasi (atau kategori
	bertriangulasi) dan $H: \mathcal{D} \to \mathcal{A}$ suatu funktor
	kohomologis. Definisikan $\mathcal{N}_H$ sebagai subkategori penuh yang
	objek-objeknya memenuhi $H(T^n X) = 0$ untuk setiap $n \in \Z$. Ini sama
	dengan mengambil $\mathcal{T}$ sebagai subkategori yang terdiri atas objek
	nol dalam Proposisi~\ref{prop:cut-subcat-by-cohomologies}. Dengan demikian,
	$\mathcal{N}_H$ merupakan subkategori prabertriangulasi (atau subkategori
	bertriangulasi).
\end{contoh}

% segment-id: o014.li2.chapter4.triangulated-basics.p008
Hasil teknis berikut akan digunakan kemudian dalam
\S~\sourcecrossref{sec:triangulated-localization}{4.3}; hasil ini bergantung
pada aksioma oktahedral (TR5).

% segment-id: o014.li2.chapter4.triangulated-basics.le002
\begin{lema}[J.-L.\ Verdier]\label{prop:triangulated-9-diag}
	Setiap diagram komutatif dalam suatu kategori bertriangulasi $\mathcal{D}$
% segment-id: o014.li2.chapter4.triangulated-basics.g015
	\[\begin{tikzcd}
		X \arrow[r, "f"] & Y \\
		X' \arrow[r, "{f'}"'] \arrow[u, "u"] & Y' \arrow[u, "v"']
	\end{tikzcd}\]
	dapat diperluas menjadi diagram
% segment-id: o014.li2.chapter4.triangulated-basics.g016
	\[\begin{tikzcd}
		TX' \arrow[r, "{Tf'}"] & TY' \arrow[r, "{Tg'}"] & TZ' \arrow[r, "{- Th'}"] & T^2 X' \arrow[phantom, ld, "\star" description] \\
		X'' \arrow[u, "o"] \arrow[r, "{f''}"] & Y'' \arrow[u, "p"] \arrow[r, "{g''}"] & Z'' \arrow[u, "q"] \arrow[r, "h''"] & TX'' \arrow[u, "{-To}"'] \\
		X \arrow[r, "f"] \arrow[u, "r"] & Y \arrow[r, "g"] \arrow[u, "s"] & Z \arrow[r, "h"] \arrow[u, "t"] & TX \arrow[u, "Tr"'] \\
		X' \arrow[r, "{f'}"'] \arrow[u, "u"] & Y' \arrow[r, "{g'}"'] \arrow[u, "v"] & Z' \arrow[r, "{h'}"'] \arrow[u, "w"] & TX' \arrow[u, "Tu"']
	\end{tikzcd}\]
	sedemikian sehingga setiap baris dan kolom merupakan segitiga terbedakan,
	persegi bertanda $\star$ antikomutatif (yakni
	$(To) h'' = -(Th') q$), dan semua persegi lainnya komutatif.
\end{lema}

% segment-id: o014.li2.chapter4.triangulated-basics.pf012
\begin{bukti}
	Kita akan membangun segitiga-segitiga terbedakan yang diminta secara
	bertahap, lalu membuktikan sifat komutatif dan antikomutatifnya. Pertama,
	terapkan (TR2) pada $u$, $v$, $f$, dan $f'$ secara terpisah untuk memperoleh
	bagian diagram
% segment-id: o014.li2.chapter4.triangulated-basics.g017
	\[\begin{tikzcd}[row sep=small, column sep=small]
		\bullet & \bullet & & \\
		\bullet \arrow[u] & \bullet \arrow[u] & & \\
		\bullet \arrow[r] \arrow[u] & \bullet \arrow[r] \arrow[u] & \bullet \arrow[r] & \bullet \\
		\bullet \arrow[r] \arrow[u] & \bullet \arrow[u] \arrow[r] & \bullet \arrow[r] & \bullet
	\end{tikzcd}\]
	yang setiap baris dan kolomnya merupakan segitiga terbedakan. Berikutnya,
	terapkan (TR2) pada $fu = vf' : X' \to Y$ untuk memperoleh segitiga
	terbedakan $X' \to Y \xrightarrow{m} A \xrightarrow{n} TX'$. Konstruksi
	berikut akan didasarkan pada (TR5).

% segment-id: o014.li2.chapter4.triangulated-basics.l003
	\begin{enumerate}
% segment-id: o014.li2.chapter4.triangulated-basics.i006
		\item Terapkan (TR5) pada tiga segitiga terbedakan yang memuat $f'$,
		$v$, dan $vf'$ untuk memperoleh morfisme di antara segitiga-segitiga
		terbedakan berikut:
% segment-id: o014.li2.chapter4.triangulated-basics.g018
		\begin{equation}\label{eqn:triangulated-9-diag-aux-0}\begin{tikzcd}
			X' \arrow[r, "{f'}"] \arrow[equal, d] & Y' \arrow[r, "{g'}"] \arrow[d, "v"] & Z' \arrow[r, "{h'}"] \arrow[d, "i"] & TX' \arrow[equal, d] \\
			X' \arrow[r, "{vf'}"] \arrow[d, "{f'}"'] & Y \arrow[r, "m"] \arrow[equal, d] & A \arrow[r, "n"] \arrow[d, "j"] & TX' \arrow[d, "{Tf'}"] \\
			Y' \arrow[r, "v"] \arrow[d, "{g'}"'] & Y \arrow[r, "s"] \arrow[d, "m"] & Y'' \arrow[r, "p"] \arrow[equal, d] & TY' \arrow[d, "{Tg'}"] \\
			Z' \arrow[r, "i"'] & A \arrow[r, "j"'] & Y'' \arrow[r, "k"'] & TZ'
		\end{tikzcd}\end{equation}
		Baris terakhir adalah segitiga terbedakan yang baru dibangun.

% segment-id: o014.li2.chapter4.triangulated-basics.i007
		\item Terapkan (TR5) pada tiga segitiga terbedakan yang memuat $u$,
		$f$, dan $fu$ untuk memperoleh morfisme di antara segitiga-segitiga
		terbedakan berikut:
% segment-id: o014.li2.chapter4.triangulated-basics.g019
		\begin{equation}\label{eqn:triangulated-9-diag-aux-1}\begin{tikzcd}
			X' \arrow[r, "u"] \arrow[equal, d] & X \arrow[r, "r"] \arrow[d, "f"] & X'' \arrow[r, "o"] \arrow[d, "a"] & TX' \arrow[equal, d] \\
			X' \arrow[r, "fu"] \arrow[d, "u"'] & Y \arrow[r, "m"] \arrow[equal, d] & A \arrow[r, "n"] \arrow[d, "b"] & TX' \arrow[d, "Tu"] \\
			X \arrow[r, "f"] \arrow[d, "r"'] & Y \arrow[r, "g"] \arrow[d, "m"] & Z \arrow[r, "h"] \arrow[equal, d] & TX \arrow[d, "Tr"] \\
			X'' \arrow[r, "a"'] & A \arrow[r, "b"'] & Z \arrow[r, "c"'] & TX''
		\end{tikzcd}\end{equation}
		Baris terakhir adalah segitiga terbedakan yang baru dibangun.

% segment-id: o014.li2.chapter4.triangulated-basics.i008
		\item Terapkan (TR2) pada $f'' := ja: X'' \to Y''$ untuk memperoleh
		segitiga terbedakan
		$X'' \xrightarrow{f''} Y'' \xrightarrow{g''} Z''
		\xrightarrow{h''} TX''$.

% segment-id: o014.li2.chapter4.triangulated-basics.i009
		\item Tinjau segitiga-segitiga terbedakan berikut, yang diperoleh dari
		langkah-langkah di atas bersama rotasi (TR3):
% segment-id: o014.li2.chapter4.triangulated-basics.q008
		\begin{gather*}
			X'' \xrightarrow{a} A \xrightarrow{b} Z \xrightarrow{c} TX'', \quad A \xrightarrow{j} Y'' \xrightarrow{k} TZ' \xrightarrow{-Ti} TA, \\
			X'' \xrightarrow{f'' = ja} Y'' \xrightarrow{g''} Z'' \xrightarrow{h''} TX'' .
		\end{gather*}
		Penerapan (TR5) kemudian memberikan morfisme di antara
		segitiga-segitiga terbedakan
% segment-id: o014.li2.chapter4.triangulated-basics.g020
		\begin{equation}\label{eqn:triangulated-9-diag-aux-2}\begin{tikzcd}
			X'' \arrow[r, "a"] \arrow[equal, d] & A \arrow[r, "b"] \arrow[d, "j"] & Z \arrow[r, "c"] \arrow[d, "t"] & TX'' \arrow[equal, d] \\
			X'' \arrow[r, "{f''}"] \arrow[d, "a"'] & Y'' \arrow[r, "{g''}"] \arrow[equal, d] & Z'' \arrow[r, "{h''}"] \arrow[d, "q"] & TX'' \arrow[d, "Ta"] \\
			A \arrow[r, "j"] \arrow[d, "b"'] & Y'' \arrow[r, "k"] \arrow[d, "m"] & TZ' \arrow[r, "{-Ti}"] \arrow[equal, d] & TA \arrow[d, "Tb"] \\
			Z \arrow[r, "t"'] & Z'' \arrow[r, "q"'] & TZ' \arrow[r, "l"'] & TZ
		\end{tikzcd}\end{equation}
		Baris terakhir adalah segitiga terbedakan yang baru dibangun. Tetapkan
		$w := bi$. Komutativitas diagram
		\eqref{eqn:triangulated-9-diag-aux-2} menyiratkan $l = -Tw$; karena itu,
		rotasi memberikan segitiga terbedakan
% segment-id: o014.li2.chapter4.triangulated-basics.q009
		\[ Z' \xrightarrow{w} Z \xrightarrow{t} Z'' \xrightarrow{q} TZ' . \]
	\end{enumerate}

% segment-id: o014.li2.chapter4.triangulated-basics.p009
	Operasi-operasi di atas memberikan baris kedua dan kolom ketiga dari diagram
	yang diminta. Baris pertama yang tersisa (atau kolom keempat) diperoleh
	dengan merotasi baris keempat (atau kolom pertama) menurut (TR3), sehingga
	baris (atau kolom) tersebut juga merupakan segitiga terbedakan (lihat
	Catatan~\ref{rem:triangle-signs}).

	Sekarang kita verifikasi sifat komutatif dan antikomutatifnya. Dari $w = bi$
	dan $s = jm$ yang disiratkan oleh
	\eqref{eqn:triangulated-9-diag-aux-0}, terlihat bahwa $(u, v, w)$ dan
	$(r, s, t)$ berturut-turut merupakan komposisi
% segment-id: o014.li2.chapter4.triangulated-basics.q010
	\begin{gather*}
		(X', Y', Z') \xrightarrow[\eqref{eqn:triangulated-9-diag-aux-0}]{(\identity, v, i)} (X', Y, A) \xrightarrow[\eqref{eqn:triangulated-9-diag-aux-1}]{(u, \identity, b)} (X, Y, Z), \\
		(X, Y, Z) \xrightarrow[\eqref{eqn:triangulated-9-diag-aux-1}]{(r, m, \identity)} (X'', A, Z) \xrightarrow[\eqref{eqn:triangulated-9-diag-aux-2}]{(\identity, j, t)} (X'', Y'', Z''),
	\end{gather*}
	sehingga keduanya merupakan morfisme di antara segitiga-segitiga. Terakhir,
	tinjau morfisme $(o, p, q, -To)$ dalam diagram. Tinggal dibuktikan bahwa
	diagram berikut komutatif:
% segment-id: o014.li2.chapter4.triangulated-basics.g021
	\[\begin{tikzcd}
		TX' \arrow[r, "{Tf'}"] & TY' \arrow[r, "{Tg'}"] & TZ' \arrow[r, "{- Th'}"] & T^2 X' \\
		A \arrow[u, "n"] \arrow[r, "j"] & Y'' \arrow[r, "k"] \arrow[u, "p"] & TZ' \arrow[equal, u] \arrow[r, "-Ti"] & TA \arrow[u, "Tn"'] \\
		X'' \arrow[u, "a"] \arrow[r, "{f''}"'] \arrow[bend left=60, uu, "o"] & Y'' \arrow[r, "{g''}"'] \arrow[equal, u] & Z'' \arrow[r, "{h''}"'] \arrow[u, "q"] & TX'' \arrow[u, "Ta"'] \arrow[bend right=60, uu, "To"'] .
	\end{tikzcd}\]
	Memang, \eqref{eqn:triangulated-9-diag-aux-1} menyiratkan $o = na$, sehingga
	kedua sayap komutatif. Karena tiga persegi di sisi kanan
	\eqref{eqn:triangulated-9-diag-aux-0} komutatif, bagian atas juga
	komutatif; komutativitas bagian bawah mengikuti dari
	\eqref{eqn:triangulated-9-diag-aux-2}. Ini membuktikan pernyataan.
\end{bukti}

% segment-id: o014.li2.chapter4.triangulated-basics.p011
Kecuali Proposisi~\ref{prop:adjoint-triangulated}, semua pernyataan dalam
bagian ini juga berlaku untuk kategori $\Bbbk$-linear, dengan $\Bbbk$ suatu
gelanggang komutatif.

% segment-id: o014.li2.chapter4.triangulated-basics.r002
\begin{catatan}
	\index{objek kompak!dalam kategori bertriangulasi}
	\index{generator!dalam kategori bertriangulasi}
	\index{Teorema Keterwakilan Brown (Brown Representability Theorem)}
	Generator dan objek kompak dalam kategori bertriangulasi merupakan topik
	penting yang tidak dibahas dalam buku ini. Walaupun gagasannya serupa dengan
	versi dalam Definisi~\ref{def:generators} dan
	\sourcecrossref{def:cpt-objects}{A.2.2}, konsep
	tersebut memerlukan definisi yang berbeda dalam kategori bertriangulasi. Hal
	ini mengarah pada bentuk penting Teorema Keterwakilan Brown untuk kategori
	bertriangulasi: jika $\mathcal{D}$ merupakan kategori bertriangulasi
	``dibangkitkan secara kompak'' yang memiliki $\coprod$ kecil, dan funktor
	kohomologis $H: \mathcal{D}^{\opp} \to \cate{Ab}$ memetakan $\coprod$ kecil
	di $\mathcal{D}$ ke $\prod$ di $\cate{Ab}$, maka $H$ terwakili. Pembaca yang
	memerlukannya disarankan merujuk pada
	\cite[Tags 09SI, 09SM, 0A8E]{stacks}.

	Perhatikan kemiripan antara Teorema Keterwakilan Brown dan berbagai
	teorema keterwakilan atau teorema funktor adjoin dalam \S\ref{sec:AFT},
	\S\ref{sec:Grothendieck-cat}, dan
	\S~\sourcecrossref{sec:cpt-objects}{A.2}.
\end{catatan}
